\documentclass[a4paper,11pt]{article}
\usepackage[utf8]{inputenc}
\usepackage[T1]{fontenc}
\usepackage[french]{babel}
\usepackage[left=2cm,right=2cm,top=2cm,bottom=2cm]{geometry}
\usepackage{xcolor}
\usepackage{hyperref}
\usepackage{fontawesome5}
\usepackage{parskip}

\definecolor{primary}{RGB}{200, 30, 30}
\hypersetup{colorlinks=true, urlcolor=primary}

\begin{document}

\noindent
\begin{minipage}[t]{0.5\textwidth}
    \textbf{Lo\"ic Aron MBASSI EWOLO}\\
    \'El\`eve-Ing\'enieur Bac+5 -- Double dipl\^ome ENSEA\\[3pt]
    \faMapMarker\ Cergy, France\\
    \faPhone\ (+33) 7 59 52 33 98\\
    \faEnvelope\ \href{mailto:loicaron.mbassiewolo@ensea.fr}{loicaron.mbassiewolo@ensea.fr}
\end{minipage}
\hfill
\begin{minipage}[t]{0.4\textwidth}
    \raggedleft
    Cergy, le 21 mars 2026
\end{minipage}

\vspace{1.2cm}

\hfill
\begin{minipage}[t]{0.5\textwidth}
    \raggedleft
    \textbf{Groupe TF1}\\
    Service Recrutement\\
    Boulogne-Billancourt, France
\end{minipage}

\vspace{1cm}

\noindent\textbf{Objet :} Candidature en alternance -- Ing\'enieur DevOps H/F -- Septembre 2026

\vspace{0.8cm}

Madame, Monsieur,

La diffusion num\'erique en continu -- streaming live, replay, publicit\'e programmatique -- impose des contraintes d'infrastructure que seules des pratiques DevOps matures permettent de satisfaire : z\'ero downtime, d\'eploiement continu, monitoring en temps r\'eel. Groupe TF1, en acc\'el\'erant sa transformation Cloud sur AWS avec Kubernetes, constitue un terrain d'apprentissage et de contribution exceptionnel pour un ing\'enieur form\'e aux architectures distribu\'ees.

Mon projet GoFind (2024-2025) repr\'esente mon exp\'erience la plus compl\`ete en DevOps : conception d'une architecture microservices (Laravel, NestJS, MongoDB, RabbitMQ), orchestration Docker Compose, mise en place de pipelines CI/CD avec tests automatis\'es. Le projet Snappy (architecture r\'eactive Spring WebFlux, caching Redis, WebSockets) m'a form\'e \`a la scalabilit\'e horizontale et au monitoring de performances d'applications temps r\'eel. Ces projets d\'emontrent une ma\^itrise concr\`ete de l'IaC et de l'automatisation du d\'eploiement.

Mon stage \`a l'INRIA Saclay (avril-ao\^ut 2026) vient compl\'eter ce profil avec l'angle infrastructure ML : d\'eploiement de pipelines de Machine Learning dans des environnements Docker, gestion de la reproductibilit\'e et de la tra\c{c}abilit\'e des exp\'eriences. Cette dimension MLOps de l'ing\'enierie DevOps est de plus en plus strat\'egique pour les groupes m\'edia qui d\'eploient des syst\`emes de recommandation et de personnalisation.

Je suis disponible d\`es septembre 2026 pour cette alternance de 12 mois \`a Boulogne-Billancourt. Contribuer \`a la fiabilit\'e et la scalabilit\'e de l'infrastructure Cloud de TF1 constitue une opportunit\'e que j'aborde avec un socle technique solide. Je reste \`a votre disposition pour un entretien.

Je vous prie d'agr\'eer, Madame, Monsieur, l'expression de mes salutations distingu\'ees.

\vspace{1cm}

\noindent\textbf{Lo\"ic Aron MBASSI EWOLO}

\end{document}
